% !TEX program = lualatex
% =============================================================================
%  Службовий лист на ПОЗДОВЖНЬОМУ бланку структурного підрозділу
%  (ДСТУ 4163:2020, додаток Б), два адресати, інтервал тексту 1,5.
% =============================================================================
\documentclass[14pt,onehalf]{dstu4163}

\dstusetup{
  blank      = long,
  parent     = {Умовна асоціація закладів освіти},
  name       = {Університет «Зразок»},
  unit       = {Кафедра прикладної економіки},
  address    = {вул. Зразкова, 15, м. Умовне, 00000},
  phone      = {(0432) 00-00-00},
  email      = {econ@zrazok.example},
  code       = {00000000},
  date       = {18.11.2026},
  index      = {07-14/221},
}

% До чотирьох адресатів (п. 5.15). Юрособі без посадової особи —
% найменування у називному відмінку.
\addressee{Університет «Приклад-1»}
\addressee{Університет «Приклад-2»\\
  Економічний факультет}

\begin{document}

\heading{Про проведення науково-практичного семінару}

\salutation{Шановні колеги!}

Кафедра прикладної економіки запрошує науково-педагогічних працівників
і здобувачів вищої освіти Ваших закладів до участі в науково-практичному
семінарі «Емпіричні методи оцінювання державної політики», який відбудеться
10 грудня 2026 року о 10:00 в онлайн-форматі.

Програмою семінару передбачено доповіді щодо методів різниці різниць,
інструментальних змінних і розривного регресійного аналізу, а також
обговорення результатів досліджень ринку житла та податкової політики.

Просимо до 01.12.2026 надіслати на адресу econ@zrazok.example список
учасників із зазначенням власного імені, прізвища, посади та теми
доповіді (за наявності).

\signature{Завідувач кафедри}{Тарас \surname{Онищенко}}

\executor{Литвин Марія (0432) 00-00-01}

\end{document}
